%% GENERATED by python/paper/build.py from results/ -- do not edit by hand.
%% Rebuild:  .venv\Scripts\python.exe python\paper\build.py --only US_CRRA_OtherShocks
%% Source:   results/shocks/US_shocksCommonX.csv
\begin{table}[!htb]
\centering
\begin{threeparttable}
\caption{Does CRRA matter for French characteristics imposed on the US -- 2020}
\label{table:US:CRRA:otherShocks}
\renewcommand{\arraystretch}{1.25}
\begin{tabularx}{.9\textwidth}{YYYYY}
\toprule
\textbf{Scenario} & \textbf{CRRA} ($\rho$) & \textbf{Tax rate} & \textbf{Savings rate} & \textbf{Avg. workweek} \\
\midrule 
 & 0.5 & 14.43\% & 15.91\% & 39.39 \\
Baseline & 1.0 & 14.43\% & 15.37\% & 39.39 \\
 & 2.0 & 14.43\% & 14.85\% & 39.39 \\[.5em]\hline\\[-.75em]
 & 0.5 & 12.89\% & $+0.82$ p.p. & 42.29 \\
Income distribution & 1.0 & 13.21\% & $+0.44$ p.p. & 41.97 \\
 & 2.0 & 13.65\% & $+0.16$ p.p. & 41.70 \\[.5em]\hline\\[-.75em]
 & 0.5 & 14.41\% & 0.00 p.p. & 32.77 \\
Leisure preferences & 1.0 & 14.43\% & 0.00 p.p. & 33.24 \\
 & 2.0 & 14.43\% & 0.00 p.p. & 33.66 \\[.5em]\hline\\[-.75em]
 & 0.5 & 14.75\% & $-0.17$ p.p. & 39.30 \\
Voting & 1.0 & 15.32\% & $-0.32$ p.p. & 39.19 \\
 & 2.0 & 16.12\% & $-0.34$ p.p. & 39.11 \\[.5em]\hline\\[-.75em]
\bottomrule
\end{tabularx}
\begin{tablenotes}
\footnotesize
\item Every $\rho$ is separately calibrated: the tax rate is a target and the workweek is normalised against each $\rho$\textquotesingle s own baseline, so both are common to the baseline rows, while the baseline savings rate is a prediction ($\beta$ is identified by the interest rate) and falls with $\rho$. The savings rate is savings relative to GDP; the baseline rows report its level and every scenario row the change against the baseline at the same $\rho$, in percentage points. Common-$X$ calibration: one leisure parameter $X$ shared across income groups, with the hours unit pinned by targeting the observed average workweek, so relative hours are a prediction rather than a calibration target.
\end{tablenotes}
\end{threeparttable}
\end{table}
